% Part: applied-modal-logics
% Chapter: temporal-logic
% Section: properties-accessibility

\documentclass[../../../include/open-logic-section]{subfiles}

\begin{document}

\olfileid[hi]{aml}{tl}{acc}

\olsection{कालिक फ़्रेमों के गुण}

चूँकि हमारे कालिक प्रतिरूप संबंध~$\prec$ पर कोई शर्त नहीं लगाते, इसलिए हमारे
परिचित अभिगृहीतों में केवल~$K$, अथवा उसके अनुरूप $K_{\Gtemp}$ और
~$K_{\Htemp}$, ही सभी प्रतिरूपों में मान्य हैं:
\begin{align*}
\tag{$K_{\Gtemp}$}  \Gtemp (p \to q) & \to (\Gtemp p \to \Gtemp q)\\
\tag{$K_{\Htemp}$}  \Htemp (p \to q) & \to (\Htemp p \to \Htemp q)
\end{align*}

परंतु यदि हम चाहते हैं कि हमारे प्रतिरूप समय के निरूपण पर अधिक कठोर शर्तें
लगाएँ---मसलन यह सुनिश्चित करके कि $\prec$ एक रैखिक क्रम हो---तो हमारे
प्रतिरूपों में अन्य सूत्र भी मान्य होंगे।

\begin{table}[t]
  \begin{tabular}{| p{.48\textwidth} || p{.48\textwidth} |}
    \hline
    {\emph{यदि $\prec$ \dots है}} & {\emph{तो \dots~$\mModel{M}$ में सत्य है:}} \\
    \hline \hline
    \emph{संक्रमणशील}: \newline
    $\forall u \forall v \forall w ((u \prec v \land v \prec w) \lif u \prec w)$ & 
    $\Ftemp \Ftemp p \lif \Ftemp p$  \\
    \hline 
    \emph{रैखिक}: \newline
    $\forall w \forall v (w \prec v \lor w = v \lor v \prec w)$ &  
    $(\Ftemp \Ptemp  p \lor \Ptemp \Ftemp p) \lif (\Ptemp p \lor p \lor \Ftemp p)$\\
    \hline
    \emph{सघन}: \newline
    $\forall w \forall v (w \prec v \to \exists u(w \prec u \land u \prec v))$ &  
    $\Ftemp p \lif \Ftemp \Ftemp p$ \\
    \hline
    \emph{अपरिबद्ध (अतीत की ओर)}: \newline
    $\forall w \exists v( v \prec w)$ &  
    $\Htemp p \to \Ptemp p$ \\
    \hline
    \emph{अपरिबद्ध (भविष्य की ओर)}: \newline
    $\forall w \exists v( w \prec v)$ &  
    $\Gtemp p \to \Ftemp p$ \\
    \hline
  \end{tabular}
  \caption{कुछ कालिक फ़्रेम-अनुरूपता गुण।}
  \ollabel{tab:correspondence}
\end{table} 

\olref{tab:correspondence} में दिए गए कई गुण समय का निरूपण करने वाले प्रतिरूप
के लिए वांछनीय विशेषताएँ लग सकते हैं। किंतु ध्यान देना चाहिए कि यद्यपि हम
$\prec$ संबंध पर अपनी इच्छानुसार शर्तें लगा सकते हैं, फिर भी प्रत्येक शर्त
ऐसे !!{formula} के अनुरूप नहीं होती जिसे कालिक तर्कशास्त्र की भाषा में व्यक्त
किया जा सके। उदाहरण के लिए, अप्रतिवर्तिता---अर्थात
$\forall w \lnot (w \prec w)$---के अनुरूप कालिक तर्कशास्त्र में कोई सूत्र
नहीं है।

\end{document}
